% =============================================================================
%  Notas del profesor — Sesión 13: Torres de Hanoi — Solución Recursiva
%  Programación Avanzada — Universidad Panamericana
%  Compilar: pdflatex sesion13_hanoi_recursivo_profesor.tex
% =============================================================================
\documentclass[12pt, oneside]{article}
\usepackage[profesor]{progavanzada}

\begin{document}

\portada{Sesión 13}{Torres de Hanoi: Solución Recursiva}{2026}

\tableofcontents
\newpage

% =============================================================================
\section{Estructura de la sesión}
% =============================================================================

\begin{center}
\begin{tabular}{clc}
  \toprule
  \textbf{Tiempo} & \textbf{Bloque} & \textbf{Tipo} \\
  \midrule
  15 min & Concurso de diseño: presentaciones + votación & Grupal \\
   5 min & Proyectar el código ganador; análisis rápido & Exposición \\
  10 min & Presentación del algoritmo recursivo & Pizarrón \\
  40 min & Actividad individual: integrar la función recursiva & Trabajo \\
   5 min & Cierre: ejecutar el programa, contar movimientos & Demo \\
   5 min & Reto extra (Fibonacci / razón áurea) para quien termina antes & Opcional \\
  \midrule
  \textbf{80 min} & & \\
  \bottomrule
\end{tabular}
\end{center}

\begin{notaprofesor}[Objetivo pedagógico de la sesión]
  El impacto de esta sesión depende del contraste con la sesión 12.
  Los alumnos pasaron 50 minutos moviendo discos a mano; hoy verán que
  \textbf{seis líneas de código} hacen exactamente lo mismo.  Ese
  contraste entre esfuerzo humano y elegancia del algoritmo recursivo
  es el momento de la sesión.  No lo adelantes: espera a que el código
  esté en pantalla y la sala vea la ejecución en vivo.
\end{notaprofesor}

% =============================================================================
\section{Cómo conducir el concurso de votación}
% =============================================================================

\subsection{Logística}

\begin{enumerate}
  \item Antes de que lleguen los alumnos, asegúrate de tener acceso a los
        archivos \texttt{.cpp} subidos a Canvas.  Si algún equipo no subió,
        acepta una foto del código en pantalla como evidencia suficiente.

  \item Cada equipo tiene \textbf{30--60 segundos} para mostrar su programa
        en pantalla (o en proyector).  No hace falta que expliquen el código;
        solo que demuestren que funciona en al menos dos estados distintos.

  \item Votación: cada persona escribe en un papel tres números de equipo
        (el suyo propio no cuenta).  Recoger los papeles o usar una votación
        de mano alzada por ronda.

  \item En caso de empate, el desempate es: el equipo cuyo código sea
        \textit{más fácil de entender en 30 segundos}.  Tú decides.
\end{enumerate}

\begin{notaprofesor}[Si el código ganador es difícil de extender]
  El código ganador puede tener el estado de las torres en variables
  globales, en un struct, con arreglos en lugar de vectores, o con
  cualquier otra estructura.  No importa: el objetivo es que el equipo
  ganador \textit{explique} su diseño al grupo en 2--3 minutos, y luego
  entre todos integren la función recursiva.

  Si el código tiene un problema estructural que haría imposible integrar
  la función recursiva (e.g., no hay forma de saber qué disco está en el
  tope de cada torre), tómalo como oportunidad didáctica: habla con el
  grupo de por qué la estructura de datos importa para la extensibilidad.
  Luego usa el segundo lugar o el código base del cuaderno de la sesión 12.
\end{notaprofesor}

% =============================================================================
\section{Presentación del algoritmo recursivo}
% =============================================================================

\subsection{Cómo presentarlo en el pizarrón}

\textbf{No des el código de inmediato.}  Primero haz las tres preguntas
siguientes y escribe las respuestas en el pizarrón:

\begin{notaprofesor}[Preguntas de arranque]
  \begin{enumerate}
    \item \textbf{``Para mover $n$ discos de A a C, ¿cuál es el último
          disco que debes mover y cuándo?''}\\
          \textit{Respuesta:} el disco $n$ (el más grande), de A a C,
          y \textit{solo puedes moverlo cuando los otros $n-1$ ya están en B}.

    \item \textbf{``Para poner los $n-1$ discos en B antes de mover el
          grande, ¿qué problema tienes que resolver?''}\\
          \textit{Respuesta:} mover $n-1$ discos de A a B usando C como
          auxiliar.  Mismo problema, $n-1$ discos.

    \item \textbf{``¿Y después de mover el disco grande?''}\\
          \textit{Respuesta:} mover $n-1$ discos de B a C usando A
          como auxiliar.  De nuevo el mismo problema.
  \end{enumerate}

  Cuando los alumnos lleguen solos a ``el mismo problema con $n-1$'',
  escribe el código en el pizarrón y deja que el silencio haga su trabajo.
\end{notaprofesor}

\subsection{El algoritmo}

La función completa:

\begin{verbatim}
void hanoi(int n, char origen, char destino, char aux,
           vector<int>& torA, vector<int>& torB, vector<int>& torC) {
    if (n == 0) return;
    hanoi(n-1, origen, aux, destino, torA, torB, torC);
    moverDisco(origen, destino, torA, torB, torC);
    hanoi(n-1, aux, destino, origen, torA, torB, torC);
}
\end{verbatim}

Si el código base usa nombres distintos de funciones o estructuras
de datos diferentes, el esquema es el mismo: el caso base devuelve
inmediatamente, y el caso general son dos sub-llamadas con el origen
y destino permutados.

\begin{notaprofesor}[Variante sin referencias a las torres]
  Si prefieres mostrar primero la versión mínima (solo imprime movimientos)
  y luego migrar al código base:

\begin{verbatim}
void hanoi(int n, char origen, char destino, char aux) {
    if (n == 0) return;
    hanoi(n-1, origen, aux, destino);
    cout << "Mueve disco de " << origen << " a " << destino << endl;
    hanoi(n-1, aux, destino, origen);
}
\end{verbatim}

  Esta versión es útil para demostrar el algoritmo en el pizarrón sin
  depender del código base del concurso.  Después los alumnos reemplazan
  el \texttt{cout} por la llamada a \texttt{moverDisco} con las torres reales.
\end{notaprofesor}

% =============================================================================
\section{El árbol de recursión — cómo explicarlo}
% =============================================================================

Dibujar el árbol en el pizarrón es más efectivo que explicarlo de palabra.
Hazlo para $n = 2$ primero (3 movimientos, árbol pequeño):

\begin{verbatim}
         hanoi(2, A, C, B)
         /       |       \
  hanoi(1,A,B,C) A→C  hanoi(1,B,C,A)
   /    |    \             /    |    \
h(0) A→B   h(0)         h(0) B→C   h(0)
\end{verbatim}

Señalar:
\begin{itemize}
  \item Cada hoja es un caso base (\texttt{n=0}, no hace nada).
  \item Cada nodo intermedio hace un movimiento (\texttt{moverDisco}).
  \item Para $n=2$: 3 nodos intermedios = 3 movimientos = $2^2 - 1$.
\end{itemize}

Luego preguntar: ``¿cuántos nodos intermedios hay para $n=3$?''
Los alumnos deben llegar a $2 \times 3 + 1 = 7 = 2^3 - 1$.

\begin{notaprofesor}[Conexión con la sesión 11]
  Esta es la recurrencia exacta que vieron en la sesión 12:
  $M(n) = 2\,M(n-1) + 1$.  Hoy la ven en el árbol de llamadas: cada
  nodo de nivel $k$ genera dos nodos de nivel $k+1$ más un movimiento.
  Es la misma ecuación.  Subrayarlo: el código y la matemática dicen
  exactamente lo mismo.
\end{notaprofesor}

% =============================================================================
\section{FAQs de los alumnos}
% =============================================================================

\begin{notaprofesor}[FAQ 1 — ¿Cómo sé qué \texttt{aux} usar en cada llamada?]
  \textbf{Pregunta:} Me confundo con qué torre poner como auxiliar en cada
  sub-llamada.\\
  \textbf{Respuesta:} Hay un truco para no confundirse.  Si quieres mover
  $n-1$ discos de \textit{origen} a \textit{aux}, entonces \textit{destino}
  queda libre y se convierte en el nuevo auxiliar.  Y viceversa.
  Resumido: lo que era \texttt{destino} pasa a ser \texttt{aux} en la primera
  sub-llamada; lo que era \texttt{origen} pasa a ser \texttt{aux} en la
  segunda.  El patrón es:
  \begin{center}
  \begin{tabular}{lll}
    \toprule
    & \textbf{origen} & \textbf{destino} \\
    \midrule
    Primera sub-llamada  & \texttt{origen} & \texttt{aux} \\
    Movimiento principal & \texttt{origen} & \texttt{destino} \\
    Segunda sub-llamada  & \texttt{aux}    & \texttt{destino} \\
    \bottomrule
  \end{tabular}
  \end{center}
\end{notaprofesor}

\begin{notaprofesor}[FAQ 2 — ¿Por qué el caso base es $n=0$ y no $n=1$?]
  \textbf{Pregunta:} ¿No sería más natural detenerse cuando queda 1 disco?\\
  \textbf{Respuesta:} Ambas versiones funcionan.  Con $n=0$ el código es
  más simétrico: cuando ya no hay discos que mover, simplemente regresa.
  Con $n=1$ como caso base, habría que escribir el movimiento de ese disco
  explícitamente.  La versión con $n=0$ deja que el caso recursivo maneje
  incluso el disco unitario, lo que hace el código más corto.
  En la práctica, muchos textos usan $n=1$; ninguna es incorrecta.
\end{notaprofesor}

\begin{notaprofesor}[FAQ 3 — ¿Cuántas llamadas hay en total?]
  \textbf{Pregunta:} Si el árbol tiene $2^n - 1$ nodos internos, ¿cuántos
  nodos totales tiene (incluyendo hojas)?\\
  \textbf{Respuesta:} Un árbol binario completo de profundidad $n$ tiene
  $2^n$ hojas y $2^n - 1$ nodos internos, en total $2^{n+1} - 1$ nodos.
  Para $n=3$: $2^4 - 1 = 15$ llamadas totales (8 de ellas con $n=0$,
  que regresan inmediatamente).
\end{notaprofesor}

\begin{notaprofesor}[FAQ 4 — ¿Se puede resolver Hanoi sin recursión?]
  \textbf{Pregunta:} ¿Existe una versión iterativa de Hanoi?\\
  \textbf{Respuesta:} Sí, pero no es trivial.  Hay dos enfoques:
  \begin{enumerate}
    \item \textbf{Simulación de stack:} usar un \texttt{stack} explícito
          que imite lo que el call stack haría de forma automática.
          Técnicamente es iterativo, pero conceptualmente es recursión disfrazada.
    \item \textbf{Regla del bit menos significativo:} en el paso $k$, mover
          el disco cuyo número corresponde al bit menos significativo de $k$
          en binario.  Esta regla (de la tabla de la sesión 12) genera la
          solución óptima sin recursión ni stack.  Es elegante pero no intuitiva.
  \end{enumerate}
  Para esta sesión no es necesario implementarlo.
\end{notaprofesor}

\begin{notaprofesor}[FAQ 5 — La función es muy lenta para $n$ grande]
  \textbf{Pregunta:} Con $n = 25$ el programa tarda mucho.  ¿Es normal?\\
  \textbf{Respuesta:} Completamente normal.  $2^{25} - 1 \approx 33$ millones
  de movimientos.  Si además dibujas las torres en cada paso, la salida en
  terminal se convierte en el cuello de botella.  La solución es no dibujar
  para $n$ grande: solo contar los movimientos.  La complejidad es $O(2^n)$
  y no hay forma de mejorarlo para el problema estándar de tres agujas.
\end{notaprofesor}

% =============================================================================
\section{Soluciones a los ejercicios}
% =============================================================================

\begin{notaprofesor}[Ejercicio 1 — Árbol de recursión para $n=2$]
  \texttt{hanoi(2, A, C, B)} genera:
  \begin{enumerate}
    \item \texttt{hanoi(1, A, B, C)} → mueve disco 1: A → B
    \item \texttt{moverDisco(A, C)} → mueve disco 2: A → C
    \item \texttt{hanoi(1, B, C, A)} → mueve disco 1: B → C
  \end{enumerate}
  Orden de movimientos: disco 1 (A→B), disco 2 (A→C), disco 1 (B→C).
  Total: 3 movimientos = $2^2 - 1$.
\end{notaprofesor}

\begin{notaprofesor}[Ejercicio 2 — Llamadas totales para $n=4$]
  Nodos internos: $2^4 - 1 = 15$ (movimientos reales).
  Hojas (\texttt{n=0}): $2^4 = 16$.
  Total: $15 + 16 = 31 = 2^5 - 1$ llamadas a la función.
  En general para $n$ discos: $2^{n+1} - 1$ llamadas totales.
\end{notaprofesor}

\begin{notaprofesor}[Ejercicio 3 — \texttt{hanoi} con indentación]
  Versión con traza similar a \texttt{sumaTraza} de la sesión 11:
\begin{verbatim}
void hanoi(int n, char origen, char destino, char aux,
           vector<int>& torA, vector<int>& torB, vector<int>& torC,
           int nivel = 0) {
    if (n == 0) return;
    string indent(nivel * 2, ' ');
    cout << indent << ">> mover " << n << " discos de "
         << origen << " a " << destino << endl;
    hanoi(n-1, origen, aux, destino, torA, torB, torC, nivel + 1);
    moverDisco(origen, destino, torA, torB, torC);
    cout << indent << "  [disco movido: " << origen
         << " → " << destino << "]" << endl;
    hanoi(n-1, aux, destino, origen, torA, torB, torC, nivel + 1);
}
\end{verbatim}
\end{notaprofesor}

\begin{notaprofesor}[Ejercicio 4 — ¿Es posible Hanoi iterativo?]
  Sí (ver FAQ 4 de esta sección).  Una respuesta completa menciona:
  la versión con stack explícito existe y es correcta; la versión basada
  en el bit menos significativo es $O(2^n)$ en tiempo pero $O(1)$ en
  espacio adicional.  Ambas son válidas.
\end{notaprofesor}

\begin{notaprofesor}[Ejercicio 5 — Parámetro $n$ con límite de visualización]
  Ejemplo de lógica en \texttt{main}:
\begin{verbatim}
int n;
cout << "¿Cuántos discos? ";
cin >> n;
if (n < 1 || n > 30) {
    cout << "Número de discos fuera de rango (1-30)." << endl;
    return 1;
}
// inicializar torres con n discos en A
// ...
bool mostrar = (n <= 20);
hanoi(n, 'A', 'C', 'B', torA, torB, torC, mostrar);
cout << "Movimientos: " << movimientos
     << " (esperado: " << (1LL << n) - 1 << ")" << endl;
\end{verbatim}
  El parámetro \texttt{mostrar} se pasa a la función y controla si
  llama a \texttt{dibujarTorres} después de cada movimiento.
\end{notaprofesor}

% =============================================================================
\section{Solución al reto extra — Fibonacci y la razón áurea}
% =============================================================================

\begin{notaprofesor}[Reto — Parte 1 y 2: implementación y tabla]
  Implementación mínima correcta:

\begin{verbatim}
long long fib(int n) {
    if (n <= 1) return n;
    return fib(n - 1) + fib(n - 2);
}

int main() {
    cout << fixed << setprecision(6);
    cout << " n     F(n)     F(n)/F(n-1)\n";
    for (int n = 2; n <= 20; n++) {
        long long fn  = fib(n);
        long long fn1 = fib(n - 1);
        cout << setw(2) << n
             << setw(9) << fn
             << "     " << (double)fn / fn1 << "\n";
    }
}
\end{verbatim}

  La razón converge a $\varphi = (1+\sqrt{5})/2 \approx 1.618034$.
  Para $n = 20$ ya tiene 6 cifras correctas.
\end{notaprofesor}

\begin{notaprofesor}[Reto — Parte 3: comparación de árboles]
  El árbol de \texttt{fib(5)}:

\begin{verbatim}
                    fib(5)
                   /      \
              fib(4)       fib(3)
             /     \       /    \
         fib(3)  fib(2) fib(2) fib(1)
         /   \   /   \  /   \
      fib(2) f(1) f(1) f(0) f(1) f(0)
      /   \
   fib(1) fib(0)
\end{verbatim}

  \textbf{Diferencia clave:} \texttt{fib(3)} aparece \textit{dos veces},
  \texttt{fib(2)} aparece \textit{tres veces}.  Hanoi, en cambio, no
  repite ningún subproblema: cada par (n, origen, destino) es único.

  Esta distinción (subproblemas superpuestos vs.\ subproblemas disjuntos)
  es el criterio que separa los problemas que se benefician de memoización
  o programación dinámica de los que no.  Es un buen gancho para el tema
  siguiente si el programa lo cubre.

  \textbf{Número de llamadas a \texttt{fib(n)}:}
  $2\,F(n+1) - 1$ en la versión recursiva ingenua.
  Para $n = 20$: $2 \times 10946 - 1 = 21891$ llamadas,
  versus solo $2^{21} - 1 = 2097151$ para Hanoi — aquí Hanoi
  es \textit{más lento}, lo que sorprende a los alumnos porque
  Fibonacci ``se ve más difícil''.  Usar ese contraste en la discusión.
\end{notaprofesor}

\begin{notaprofesor}[Cómo presentar el reto en clase]
  No explicar el reto al grupo: entregarlo impreso o proyectarlo y dejarlo
  trabajar en silencio.  Quien termina la actividad principal con tiempo
  de sobra avanza solo.

  Si al final hay tiempo de discusión, la pregunta que más reacción genera
  es: ``¿Cuántas veces se llama \texttt{fib(20)}?''  La mayoría esperará
  un número cercano a 20; la respuesta correcta ($21\,891$) los sorprende
  y abre la conversación sobre memoización.
\end{notaprofesor}

% =============================================================================
\section{Cierre del arco de tres sesiones}
% =============================================================================

Al final de la sesión, vale la pena hacer una síntesis breve en pizarrón:

\begin{center}
\begin{tabular}{cll}
  \toprule
  \textbf{Sesión} & \textbf{Tema} & \textbf{Qué aprendieron} \\
  \midrule
  11 & Funciones recursivas   & Caso base, caso recursivo, call stack \\
  12 & Hanoi a mano + diseño  & El problema ``se siente'' antes de resolverse \\
  13 & Algoritmo recursivo    & 6 líneas hacen lo que tomó 50 min a mano \\
  \bottomrule
\end{tabular}
\end{center}

\begin{notaprofesor}[Frase de cierre sugerida]
  ``La recursión no es un truco del programador: es la forma natural de
  describir un problema que se repite a menor escala.  Hanoi lo hace
  evidente porque la solución recursiva \textit{es} la definición del
  problema, escrita en C++.''
\end{notaprofesor}

% =============================================================================
\section{Resumen para el profesor}
% =============================================================================

\begin{recuerda}
\begin{itemize}
  \item El concurso toma 15 minutos; tenerlo preparado antes de que lleguen
        los alumnos acelera el proceso.
  \item No dar el código antes de las tres preguntas de arranque.
  \item La actividad individual dura 45 minutos; circular para verificar que
        todos hayan integrado \texttt{moverDisco} (no solo impriman movimientos).
  \item Cerrar con la ejecución en vivo y el conteo de movimientos.
  \item Al final del arco de tres sesiones, conectar con el siguiente tema:
        las funciones recursivas son la base de algoritmos de ordenamiento
        (mergesort, quicksort) que verán más adelante.
\end{itemize}
\end{recuerda}

\end{document}
